\documentclass[aspectratio=169, 11pt, xcolor={table}]{beamer}
% Shared style for the algorithmics course slides (same template as the
% ENSB generative-AI course), with C code listings added.
% Dark blue removed: every slide (cover included) uses the light background
% of the template's content slides, with teal + orange as accents.

\usepackage{amsmath}
\usepackage{amssymb}
\usepackage{booktabs}
\usepackage{colortbl}
\usepackage{tabularx}
\usepackage{tikz}
\usepackage[skins]{tcolorbox}
\usepackage{fontawesome5}
\usepackage{listings}
\usepackage{polyglossia}

\usetikzlibrary{positioning, calc, arrows.meta, shapes.geometric}

%------------------------------------------------
% Languages and fonts (compile from slides/ with XeLaTeX or tectonic)
%------------------------------------------------
\setdefaultlanguage[numerals=maghrib]{arabic}
\setotherlanguage{english}
\usefonttheme{serif}

\newfontfamily\arabicfont[
    Script=Arabic, Scale=0.98, Path=../../fonts/, Extension=.ttf,
    UprightFont=*-Regular, BoldFont=*-Bold,
    ItalicFont=*-Italic, SlantedFont=*-Italic,
    BoldItalicFont=*-BoldItalic, BoldSlantedFont=*-BoldItalic]{Amiri}
\newfontfamily\arabicfontsf[
    Script=Arabic, Scale=0.98, Path=../../fonts/, Extension=.ttf,
    UprightFont=*-Regular, BoldFont=*-Bold,
    ItalicFont=*-Italic, SlantedFont=*-Italic,
    BoldItalicFont=*-BoldItalic, BoldSlantedFont=*-BoldItalic]{Amiri}
% Arabic fallback for \ttfamily (needed by the C listings)
\newfontfamily\arabicfonttt[Script=Arabic, Scale=0.98, Path=../../fonts/,
    Extension=.ttf, UprightFont=*-Regular, BoldFont=*-Bold,
    ItalicFont=*-Italic, SlantedFont=*-Italic,
    BoldItalicFont=*-BoldItalic, BoldSlantedFont=*-BoldItalic]{Amiri}
% Latin monospace font for C code (selected directly, so polyglossia
% does not swap it for the Arabic fallback above)
\newfontfamily\codefont[Scale=0.95, BoldFont=lmmonolt10-bold.otf,
    ItalicFont=lmmono10-italic.otf]{lmmono10-regular.otf}
% Heading font
\newfontfamily\khatt[Script=Arabic, Scale=1.05, Path=../../fonts/,
    BoldFont=ae_AlBattar.ttf]{ae_AlBattar.ttf}

%------------------------------------------------
% Colors (template palette without the navy tones)
%------------------------------------------------
\definecolor{brandTeal}{HTML}{5CC8E6}
\definecolor{brandOrange}{HTML}{FFA27F}
\definecolor{brandLight}{HTML}{F8F9FA}
\definecolor{brandDark}{HTML}{263238}
\definecolor{brandMuted}{HTML}{6C757D}
\colorlet{tealDeep}{brandTeal!60!black}
\definecolor{codeGrey}{HTML}{F1F3F5}

\setbeamercolor{background canvas}{bg=brandLight}
\setbeamercolor{normal text}{fg=brandDark}
\setbeamercolor{structure}{fg=brandTeal}
\setbeamertemplate{navigation symbols}{}

% beamer.cls sets \raggedright globally; under RTL that centers short lines
% or pushes them left, so go back to justified text (flush right).
\rightskip=0pt\relax
% Lists call \raggedright too; make it a no-op for RTL slides.
\AtBeginDocument{\renewcommand{\raggedright}{\rightskip=0pt\relax}}

%------------------------------------------------
% Frame title, footer, lists
%------------------------------------------------
\setbeamertemplate{frametitle}{%
    \vspace{0.35cm}
    \noindent{\khatt\bfseries\Large\color{brandDark}\insertframetitle}
    \par
    \vskip-2pt
    \noindent{\color{brandTeal}\rule{\textwidth}{1.5pt}}
    \vspace{0.05cm}
}

% Footer: page number only
\setbeamertemplate{footline}{%
    \hbox to \paperwidth{\hfill
        {\color{brandMuted}\scriptsize\LR{\insertframenumber\,/\,\inserttotalframenumber}}%
        \hspace{0.6cm}}%
    \vspace{0.25cm}}

\setbeamertemplate{itemize item}{\color{brandTeal}\small\faChevronLeft}
\setbeamertemplate{itemize subitem}{\color{brandMuted}\small\faAngleLeft}
\setbeamertemplate{enumerate item}{\color{tealDeep}\bfseries\insertenumlabel)}

% Tables: every cell framed by a visible border
\newcolumntype{R}{>{\raggedleft\arraybackslash}X}
\newcolumntype{C}{>{\centering\arraybackslash}X}
\newcommand{\ras}[1]{\cellcolor{brandTeal!20}\textbf{#1}}
% Colored text inside a table cell (\leavevmode keeps it on the row baseline)
\newcommand{\cell}[2]{\leavevmode\textcolor{#1}{#2}}
\renewcommand{\arraystretch}{1.35}
\setlength{\arrayrulewidth}{0.6pt}
\arrayrulecolor{brandMuted!70}

%------------------------------------------------
% Boxes (the two boxes of the template's 2nd slide)
%------------------------------------------------
\tcbset{
    brandbox/.style={
        enhanced, arc=4pt, outer arc=4pt, boxrule=1.5pt,
        coltitle=white, coltext=brandDark, fonttitle=\khatt\bfseries,
        toptitle=0.8mm, bottomtitle=0.8mm,
        top=1.2mm, bottom=1.2mm, left=3mm, right=3mm,
        before skip=4pt, after skip=4pt}
}
% Teal box: definitions, objectives, summaries
\newtcolorbox{tbox}[2][]{brandbox, colback=brandTeal!10,
    colframe=brandTeal, title={#2}, #1}
% Orange box: warnings
\newtcolorbox{obox}[2][]{brandbox, colback=brandOrange!10,
    colframe=brandOrange, title={#2}, #1}
% Card (template's 3rd slide, in teal)
\newtcolorbox{card}[2][]{enhanced, colback=brandTeal!6, colframe=brandTeal,
    arc=6pt, boxrule=1.2pt, halign=center, halign title=center,
    coltitle=white, coltext=brandDark, fonttitle=\khatt\bfseries,
    top=1.5mm, bottom=1.5mm, left=2mm, right=2mm, title={#2}, #1}
% Prompt text written to the tool
\newtcolorbox{amr}[1][brandTeal]{enhanced, arc=1mm, boxrule=0pt,
    colback=white, overlay={\fill[#1] (frame.north east) rectangle ([xshift=-2.5pt]frame.south east);},
    left=6pt, right=8pt, top=2pt, bottom=2pt,
    before skip=3pt, after skip=3pt}

%------------------------------------------------
% TikZ (Arabic right-to-left inside nodes)
%------------------------------------------------
\newcommand{\nodecolor}{}
\tikzset{
    every node/.append style={
        execute at begin node={\nodecolor\everypar{\beginR}\beginR},
        execute at end node=\endR},
    lawn/.code={\def\nodecolor{\color{#1}}},
    khatwa/.style={draw=#1, fill=#1!12, very thick, rounded corners=5pt,
        minimum height=0.95cm, text width=1.9cm, align=center,
        font=\khatt\bfseries\small},
    khatwa/.default=brandTeal,
    sahm/.style={-{Stealth[length=2.5mm]}, thick, color=brandMuted},
    rang/.style={draw=#1, fill=#1!12},
    raqm/.style={circle, fill=#1, lawn=white, inner sep=1.5pt,
        font=\small\bfseries},
}

% Horizontal chain of steps, right to left:
% \flow[width]{{a}, {b\\c}, {d}}
\newcommand{\flow}[2][1.9cm]{%
    \begin{tikzpicture}[node distance=5mm]
        \foreach \t [count=\i] in {#2} {
            \ifnum\i=1
                \node[khatwa, text width=#1] (f\i) {\t};
            \else
                \pgfmathtruncatemacro\p{\i-1}
                \node[khatwa, text width=#1, left=of f\p] (f\i) {\t};
                \draw[sahm] (f\p) -- (f\i);
            \fi
        }
    \end{tikzpicture}}

%------------------------------------------------
% C code (always left-to-right, inside a fragile frame)
%------------------------------------------------
\lstdefinestyle{cstyle}{
    language=C, basicstyle=\codefont\footnotesize, columns=fullflexible,
    keepspaces=true, showstringspaces=false, upquote=true,
    keywordstyle=\color{tealDeep}\bfseries,
    commentstyle=\color{red!75!black},
    stringstyle=\color{brandOrange!85!black},
    xleftmargin=16pt,
    numbers=left, numberstyle=\tiny\color{brandMuted}, numbersep=7pt,
    aboveskip=0pt, belowskip=0pt,
    morekeywords={bool, true, false}}
% Inline listings must sit in the english environment, not LTR, otherwise
% numbers such as 3.14 are reversed.
% Numbers in Arabic text: \textenglish{1.75} (not \LR).
% Compact code size for long programs (fixed line height, no gaps)
\newcommand{\codesmall}{\fontsize{8}{8}\selectfont}
\newcommand{\codetiny}{\fontsize{7.4}{7.6}\selectfont}
% Inline C: \code{printf}
\newcommand{\code}[1]{\mbox{\textenglish{\codefont\begingroup\color{tealDeep}#1\endgroup}}}
% Program output (console)
\lstdefinestyle{outstyle}{
    language={}, basicstyle=\codefont\footnotesize\color{brandDark},
    columns=fullflexible, keepspaces=true, showstringspaces=false,
    aboveskip=0pt, belowskip=0pt}

% Listings read from files (no fragile frame needed):
%   \ccode{d1/circle.c}   \out{d1/circle.out}
% One box per listing for the grey background (listings' own per-line
% background leaves thin white gaps between lines)
\newtcolorbox{codebox}{enhanced, colback=codeGrey, colframe=codeGrey,
    boxrule=0pt, arc=0pt, left=4pt, right=4pt, top=1pt, bottom=1pt,
    before skip=0pt, after skip=2pt,
    overlay={\fill[brandMuted!45] (frame.south west) rectangle ([xshift=2.5pt]frame.north west);}}
\newtcolorbox{outbox}{enhanced, colback=brandMuted!20, colframe=brandMuted!20,
    boxrule=0pt, arc=0pt, left=4pt, right=4pt, top=1.5pt, bottom=1.5pt,
    before skip=2pt, after skip=2pt}
\newcommand{\ccode}[2][]{\begin{codebox}\begin{english}\lstinputlisting[style=cstyle,#1]{../code/#2}\end{english}\end{codebox}}
\newcommand{\out}[2][]{\begin{outbox}\begin{english}\lstinputlisting[style=outstyle,#1]{../code/#2}\end{english}\end{outbox}}
% Audience line of the cover
\newcommand{\audience}{طلبة السنة الأولى}

%------------------------------------------------
% Cover (template's 1st slide on the light background)
%------------------------------------------------
\newcommand{\cover}[3]{%
\begin{frame}[plain]
    \begin{tikzpicture}[remember picture, overlay]
        \draw[brandTeal, line width=1.5pt, opacity=0.3]
            (current page.south east) ++(-1,1) circle (4cm);
        \draw[brandTeal, line width=0.8pt, opacity=0.25]
            (current page.north west) ++(2,2) circle (6cm);
        \draw[brandOrange, line width=0.5pt, opacity=0.25]
            (current page.north west) ++(2,2) circle (7cm);
        \node[anchor=north east] at ($(current page.north east)+(-1.0,-0.8)$) {
            \begin{tikzpicture}[scale=0.35]
                \fill[brandTeal] (0,0) -- (0.8,1.2) -- (1.6,0) -- cycle;
                \fill[brandOrange] (0.8,1.2) -- (1.6,2.4) -- (2.4,1.2) -- cycle;
            \end{tikzpicture}
        };
    \end{tikzpicture}

    \vspace{1.6cm}
    \begin{minipage}{0.85\textwidth}
        {\color{tealDeep}\khatt\bfseries #1} \\[0.3cm]
        {\raggedleft\color{brandDark}\Huge\khatt\bfseries #2\par} \vspace{0.3cm}
        {\color{brandMuted}\large #3}
    \end{minipage}

    \vspace{1.2cm}
    \begin{minipage}{0.85\textwidth}
        \color{brandMuted}\small
        \begin{tabular}{@{}ll}
            \textbf{الأستاذ:} & د. بن عبد الرزاق يوسف \\
            \textbf{الفئة:}   & \audience \\
        \end{tabular}
    \end{minipage}
\end{frame}}

\title{الحلقات التكرارية}

\begin{document}

\cover{درس الفصل الثالث}
{الحلقات التكرارية}
{\textenglish{for} و\textenglish{while} و\textenglish{do while} بلغة \textenglish{C}}

\begin{frame}{أهداف الدرس}
    \begin{tbox}{\faBullseye\quad في نهاية هذا الدرس، تستطيع أن:}
        \begin{itemize}
            \item تكرر تعليمات باستعمال \code{for} و\code{while} و\code{do while}.
            \item تختار الحلقة المناسبة لكل مشكلة.
            \item تتتبع تنفيذ حلقة وتتجنب الحلقة اللانهائية.
        \end{itemize}
    \end{tbox}
\end{frame}

\begin{frame}{لماذا الحلقات؟}
    \begin{columns}[T]
        \column{0.45\textwidth}
        \ccode{d3/noloop.c}
        \column{0.5\textwidth}
        لعرض الأعداد من 1 إلى 5 نحتاج إلى خمس تعليمات.\\[4pt]
        ولعرض الأعداد من 1 إلى 1000؟ نستعمل \textbf{حلقة}.
    \end{columns}
\end{frame}

\begin{frame}{الحلقة}
    \begin{tbox}{\faLightbulb\quad تعريف}
        \textbf{الحلقة}: بنية تسمح بتكرار مجموعة من التعليمات
        ما دام \textbf{شرط} معين صحيحًا.
    \end{tbox}
\end{frame}

\begin{frame}{أنواع الحلقات في \textenglish{C}}
    \begin{center}
    \begin{tabular}{|c|c|}
        \hline
        \ras{الحلقة} & \ras{نستعملها عندما} \\ \hline
        \code{for}      & نعرف عدد التكرارات مسبقًا \\ \hline
        \code{while}    & لا نعرف عدد التكرارات، وقد لا تُنفذ أبدًا \\ \hline
        \code{do while} & تُنفذ مرة واحدة على الأقل \\ \hline
    \end{tabular}
    \end{center}
\end{frame}

\begin{frame}{حلقة \textenglish{for}}
    \begin{columns}[T]
        \column{0.5\textwidth}
        \ccode{d3/for.c}
        \column{0.3\textwidth}
        \out{d3/for.out}
    \end{columns}
\end{frame}

\begin{frame}{أجزاء حلقة \textenglish{for}}
    \vspace{0.6cm}
    \begin{center}
        \flow[2.6cm]{{البداية\\\code{i = 1}}, {الشرط\\\code{i <= 5}}, {التعليمات}, {الزيادة\\\code{i++}}}
    \end{center}
    \vspace{0.3cm}
    بعد الزيادة نعود إلى الشرط. تتوقف الحلقة عندما يصبح الشرط \textbf{خاطئًا}.
\end{frame}

\begin{frame}{مثال: مجموع الأعداد من 1 إلى \textenglish{n}}
    \begin{columns}[T]
        \column{0.55\textwidth}
        \ccode{d3/sum.c}
        \column{0.35\textwidth}
        \out{d3/sum.out}
    \end{columns}
\end{frame}

\begin{frame}{جدول التتبع: \textenglish{n = 4}}
    \begin{center}
    \begin{tabular}{|c|c|c|c|}
        \hline
        \ras{\code{i}} & \ras{\code{i <= n}} & \ras{\code{s = s + i}} & \ras{\code{s}} \\ \hline
        1 & صحيح & $0 + 1$ & 1 \\ \hline
        2 & صحيح & $1 + 2$ & 3 \\ \hline
        3 & صحيح & $3 + 3$ & 6 \\ \hline
        4 & صحيح & $6 + 4$ & 10 \\ \hline
        5 & \cell{brandOrange!80!black}{خطأ} & -- & \cell{tealDeep}{\textbf{10}} \\ \hline
    \end{tabular}
    \end{center}
    نعطي \code{s} القيمة الأولية \textenglish{0} قبل الحلقة.
\end{frame}

\begin{frame}{حلقة \textenglish{while}}
    \begin{columns}[T]
        \column{0.5\textwidth}
        \ccode{d3/while.c}
        \column{0.45\textwidth}
        نختبر الشرط \textbf{قبل} كل تكرار.\\[4pt]
        لا ننسى تغيير \code{i} داخل الحلقة.
    \end{columns}
\end{frame}

\begin{frame}{مثال: عدد أرقام عدد صحيح}
    \begin{columns}[T]
        \column{0.55\textwidth}
        \ccode{d3/digits.c}
        \column{0.35\textwidth}
        \out{d3/digits.out}
        \vspace{0.3cm}
        لا نعرف عدد التكرارات مسبقًا، لذلك نستعمل \code{while}.
    \end{columns}
\end{frame}

\begin{frame}{حلقة \textenglish{do while}}
    \begin{columns}[T]
        \column{0.5\textwidth}
        \ccode{d3/dowhile.c}
        \column{0.45\textwidth}
        \out{d3/dowhile.out}
    \end{columns}
    \vspace{0.3cm}
    نختبر الشرط \textbf{بعد} التكرار، إذن تُنفذ الحلقة مرة واحدة على الأقل.
\end{frame}

\begin{frame}{الحلقات المتداخلة}
    \begin{columns}[T]
        \column{0.55\textwidth}
        \ccode{d3/nested.c}
        \column{0.35\textwidth}
        \out{d3/nested.out}
    \end{columns}
    \vspace{0.2cm}
    لكل قيمة من \code{i} تُنفذ الحلقة الداخلية كاملة.
\end{frame}

\begin{frame}{\textenglish{break} و\textenglish{continue}}
    \begin{columns}[T]
        \column{0.6\textwidth}
        \ccode{d3/break.c}
        \column{0.3\textwidth}
        \out{d3/break.out}
    \end{columns}
    \vspace{0.3cm}
    \code{continue}: تنتقل إلى التكرار التالي. \code{break}: تُوقف الحلقة.
\end{frame}

\begin{frame}{مثال: العاملي \textenglish{n!}}
    \begin{columns}[T]
        \column{0.55\textwidth}
        \ccode{d3/fact.c}
        \column{0.35\textwidth}
        \out{d3/fact.out}
    \end{columns}
    \vspace{0.2cm}
    \code{5! = 1 * 2 * 3 * 4 * 5 = 120}
\end{frame}

\begin{frame}{الحلقة اللانهائية}
    \begin{columns}[T]
        \column{0.55\textwidth}
        \ccode{d3/infinite.c}
        \column{0.4\textwidth}
        \begin{obox}{\faExclamationTriangle\quad انتبه}
            الشرط يبقى صحيحًا دائمًا، فلا تتوقف الحلقة.
        \end{obox}
    \end{columns}
\end{frame}

\begin{frame}{أمثلة تطبيقية}
    \begin{tbox}{\faLaptopCode\quad عشرة أمثلة من مجالات مختلفة}
        نطبق ما تعلمناه في مجالات من الحياة اليومية:
        الصحة، الرياضة، الطقس، التجارة الإلكترونية، البيئة، الأمن السيبراني، ألعاب الفيديو، الهواتف الذكية، الفضاء، الشبكات الاجتماعية.
    \end{tbox}
\end{frame}

\begin{frame}[t]{المثال الأول: الصحة}
    \begin{columns}[T]
        \column{0.64\textwidth}
        \ccode[basicstyle=\codefont\codesmall]{d3/ex01.c}
        \column{0.33\textwidth}
        {\small\textbf{عدد الخطوات في أسبوع}: نجمع خطوات كل يوم ثم نحسب المعدل اليومي.\par}
        \vspace{0.3cm}
        \out[basicstyle=\codefont\scriptsize\color{brandDark}]{d3/ex01.out}
    \end{columns}
\end{frame}

\begin{frame}[t]{المثال الثاني: الرياضة}
    \begin{columns}[T]
        \column{0.64\textwidth}
        \ccode[basicstyle=\codefont\codesmall]{d3/ex02.c}
        \column{0.33\textwidth}
        {\small\textbf{برنامج دوري}: حلقتان متداخلتان لعرض كل المباريات بين 4 فرق.\par}
        \vspace{0.3cm}
        \out[basicstyle=\codefont\scriptsize\color{brandDark}]{d3/ex02.out}
    \end{columns}
\end{frame}

\begin{frame}[t]{المثال الثالث: الطقس}
    \begin{columns}[T]
        \column{0.64\textwidth}
        \ccode[basicstyle=\codefont\codesmall]{d3/ex03.c}
        \column{0.33\textwidth}
        {\small\textbf{جدول تحويل الحرارة}: حلقة \code{for} بخطوة 10 لعرض جدول تحويل.\par}
        \vspace{0.3cm}
        \out[basicstyle=\codefont\scriptsize\color{brandDark}]{d3/ex03.out}
    \end{columns}
\end{frame}

\begin{frame}[t]{المثال الرابع: التجارة الإلكترونية}
    \begin{columns}[T]
        \column{0.64\textwidth}
        \ccode[basicstyle=\codefont\codesmall]{d3/ex04.c}
        \column{0.33\textwidth}
        {\small\textbf{نمو المبيعات}: المبيعات تزيد 10\% كل شهر: متى نصل إلى الهدف؟\par}
        \vspace{0.3cm}
        \out[basicstyle=\codefont\scriptsize\color{brandDark}]{d3/ex04.out}
    \end{columns}
\end{frame}

\begin{frame}[t]{المثال الخامس: البيئة}
    \begin{columns}[T]
        \column{0.64\textwidth}
        \ccode[basicstyle=\codefont\codesmall]{d3/ex05.c}
        \column{0.33\textwidth}
        {\small\textbf{انخفاض التلوث}: مؤشر التلوث ينقص 15\% كل سنة حتى يصبح أقل من 50.\par}
        \vspace{0.3cm}
        \out[basicstyle=\codefont\scriptsize\color{brandDark}]{d3/ex05.out}
    \end{columns}
\end{frame}

\begin{frame}[t]{المثال السادس: الأمن السيبراني}
    \begin{columns}[T]
        \column{0.64\textwidth}
        \ccode[basicstyle=\codefont\codesmall]{d3/ex06.c}
        \column{0.33\textwidth}
        {\small\textbf{تخمين رمز سري}: نجرب كل الرموز من 0000 إلى 9999، ونتوقف بـ\code{break} عند إيجاد الرمز.\par}
        \vspace{0.3cm}
        \out[basicstyle=\codefont\scriptsize\color{brandDark}]{d3/ex06.out}
    \end{columns}
\end{frame}

\begin{frame}[t]{المثال السابع: ألعاب الفيديو}
    \begin{columns}[T]
        \column{0.64\textwidth}
        \ccode[basicstyle=\codefont\codesmall]{d3/ex07.c}
        \column{0.33\textwidth}
        {\small\textbf{مستويات اللاعب}: كل مستوى يحتاج ضعف نقاط الخبرة السابقة.\par}
        \vspace{0.3cm}
        \out[basicstyle=\codefont\scriptsize\color{brandDark}]{d3/ex07.out}
    \end{columns}
\end{frame}

\begin{frame}[t]{المثال الثامن: الهواتف الذكية}
    \begin{columns}[T]
        \column{0.64\textwidth}
        \ccode[basicstyle=\codefont\codesmall]{d3/ex08.c}
        \column{0.33\textwidth}
        {\small\textbf{تفريغ البطارية}: البطارية تفقد 7\% كل ساعة: نستعمل \code{do while}.\par}
        \vspace{0.3cm}
        \out[basicstyle=\codefont\scriptsize\color{brandDark}]{d3/ex08.out}
    \end{columns}
\end{frame}

\begin{frame}[t]{المثال التاسع: الفضاء}
    \begin{columns}[T]
        \column{0.64\textwidth}
        \ccode[basicstyle=\codefont\codesmall]{d3/ex09.c}
        \column{0.33\textwidth}
        {\small\textbf{العد التنازلي}: حلقة \code{for} تنازلية قبل إطلاق صاروخ.\par}
        \vspace{0.3cm}
        \out[basicstyle=\codefont\scriptsize\color{brandDark}]{d3/ex09.out}
    \end{columns}
\end{frame}

\begin{frame}[t]{المثال العاشر: الشبكات الاجتماعية}
    \begin{columns}[T]
        \column{0.64\textwidth}
        \ccode[basicstyle=\codefont\codesmall]{d3/ex10.c}
        \column{0.33\textwidth}
        {\small\textbf{منشور ينتشر بسرعة}: عدد المشاهدات يتضاعف 3 مرات كل ساعة: متى نصل إلى مليون؟\par}
        \vspace{0.3cm}
        \out[basicstyle=\codefont\scriptsize\color{brandDark}]{d3/ex10.out}
    \end{columns}
\end{frame}


\begin{frame}{أخطاء شائعة}
    \begin{obox}{\faExclamationTriangle\quad تجنّب هذه الأخطاء}
        \begin{itemize}
            \item وضع \code{;} بعد \code{for (...)} أو \code{while (...)}.
            \item نسيان تهيئة المجموع \code{s = 0} أو الجداء \code{f = 1}.
            \item خطأ في الحدود: \code{i < n} بدل \code{i <= n}.
        \end{itemize}
    \end{obox}
\end{frame}

\begin{frame}{خلاصة الدرس}
    \begin{tbox}{\faCheckCircle\quad خلاصة الدرس}
        \begin{itemize}
            \item \code{for}: عدد التكرارات معروف.
            \item \code{while}: الشرط قبل التكرار، و\code{do while}: الشرط بعده.
            \item يجب أن يصبح الشرط خاطئًا في لحظة ما لتتوقف الحلقة.
        \end{itemize}
    \end{tbox}
\end{frame}

\end{document}
